\section{Final Considerations}

This paper presented a thorough investigation of machine learning architectures for cross-domain image retrieval, ranging from traditional CNNs to modern foundation and transformer-based models. The most striking finding was the clear dominance of CLIP RN50x64, which achieved a weighted top-k accuracy score of 734.63 out of 1000 points, substantially outperforming all other evaluated approaches. This considerable margin reflects a fundamental shift in retrieval effectiveness driven by semantic understanding rather than purely visual similarity.

Traditional CNN architectures, even when fine-tuned on domain-specific data, struggled to bridge the natural-to-synthetic domain gap. These models primarily rely on pixel-level and texture-based features, which are brittle when faced with differences in lighting, textures, and rendering artifacts commonly found in synthetic images. This limitation was evident in their poor performance on the competition task, which required matching natural celebrity photographs with AI-generated synthetic counterparts.

CLIP RN50x64, built on a wide ResNet-50 backbone with 64 channels per layer, significantly wider than the standard ResNet-50, integrates contrastive multimodal training that aligns image and text embeddings within a shared semantic space. This enables the model to learn conceptual relationships that transcend superficial visual differences, allowing it to robustly identify the same individual in both natural and synthetic images without the need for fine-tuning on specific datasets, as demonstrated by Iijima et al.~\cite{iijima2024multimodal}. Among the CLIP variants evaluated, RN50x64 was best suited for this task due to its larger model capacity and more extensive training, which allowed it to develop richer and more discriminative semantic features.

The use of the “Garden of Eden” flower dataset for pre-competition evaluation, while providing useful insights into relative model capabilities, was limited by its lack of domain mismatch, which was a defining feature of the actual competition. The stark contrast in performance between DINO, which excelled on the homogeneous flower dataset, and CLIP, which outperformed all others on the celebrity matching task, highlights the critical role of training methodology over architectural complexity or model size. DINO, despite its transformer-based architecture and global reasoning abilities~\cite{caron2021emerging}, relies solely on self-supervised learning of natural images, which promotes appearance-level consistency rather than semantic identity preservation. This causes it to focus on mid-level or instance-agnostic patterns~\cite{amir2021deep}, making it effective for clustering similar textures and shapes but insufficient for associating different domains of the same concept, such as real and synthetic portraits of the same person.

Overall, the cross-domain nature of the competition exposed fundamental limitations of traditional computer vision approaches. The results demonstrate that pixel-level similarity, as typically captured by CNN features, is insufficient for bridging significant domain gaps. CLIP’s exceptional performance underscores the essential distinction between visual similarity and semantic similarity, suggesting that exposure to diverse image-text pairs during pre-training encourages the development of semantically meaningful, identity-preserving features that generalize robustly across domains. This makes CLIP RN50x64 especially suitable for cross-modal and cross-domain retrieval tasks.

These findings have important implications for the future development of image retrieval systems, indicating that foundation models trained on large-scale, diverse multimodal datasets may offer more robust and generalizable solutions than traditional fine-tuned domain-specific models. The results also emphasize that contrastive multimodal training can create more effective semantic representations than self-supervised visual learning alone, even when the latter employs sophisticated transformer architectures. This insight should guide future research toward multimodal and semantic-based learning approaches for retrieval and related computer vision tasks.

